\chapter{Web II: the Main-Loop Problem}
\label{ch:web-main-loop}

A native \texttt{Game::Run()} owns a straightforward loop: construct the game, enter the loop,
return when it ends, and then destroy the object.  Emscripten inverts that control.  The program
registers a callback and yields to JavaScript so the browser can schedule frames.  CNA once
treated that yield like an ordinary returning call, creating a lifetime bug before the first
frame was drawn.

\section{The throw that is not a C\texttt{++} exception}

The Web path calls \texttt{emscripten\_set\_main\_loop} with
\texttt{simulateInfiniteLoop=1}.  Emscripten implements the apparent non-return by emitting a
raw JavaScript \texttt{throw 'unwind'}.  With WebAssembly exception handling enabled by
\texttt{-fwasm-exceptions}, CNA's generated C\texttt{++} cleanup path intercepted that foreign
throw and unwound the current stack.

That interaction made this apparently conventional entry point unsafe:

\begin{verbatim}
int main() {
    MyGame game;
    game.Run();
}
\end{verbatim}

The unwind destroyed \texttt{game}, while the registered callback still retained pointers into
it, including the graphics-device manager.  The page could build and start successfully, yet
the first scheduled frame dereferenced dead state.  This is a useful counterexample to the
idea that compilation plus module startup proves a viable Web port.

\section{Reducing the failure before explaining it}

The first hypothesis blamed multiple inheritance and pointer adjustment in the graphics
manager.  It was plausible: a stale-looking address crossed an interface boundary, and Web
builds expose layouts differently from familiar native debuggers.  Focused reproductions under
AddressSanitizer, UndefinedBehaviorSanitizer, vptr checks, and Emscripten \texttt{SAFE\_HEAP}
did not support it.

An isolated lifetime reproduction did.  Destruction occurred during the control hand-off,
before frame one, and the callback later observed the destroyed object.  The decisive evidence
was temporal rather than structural: constructor, hand-off, destructor, callback.  Recording
that sequence ruled out a renderer-specific theory and located the defect in the hosting
contract shared by every Web renderer.

\section{The repair and its boundary}

The repaired examples allocate the game for browser lifetime rather than automatic scope:

\begin{verbatim}
int main() {
    auto* game = new MyGame();
    game->Run();
}
\end{verbatim}

Twelve example entry points were changed in commit \texttt{b6686bfba}.  Intentional process-
lifetime ownership is appropriate here because the browser loop does not return to a normal
destruction point.  A richer host could instead hold the object in explicitly managed state
and delete it after cancelling the loop, but stack lifetime cannot cross this hand-off.

The repair was not mechanically complete.  The renderer benchmark still constructs its game
on the stack at \texttt{modules/graphics/examples/graphics\_renderer\_benchmark.cpp:274}.  It
is an Emscripten-only target and remains a regression-shaped exception.  This matters as much
as the twelve repaired sites: a search-based migration must end with a search for survivors,
not with a count of edited files.

\section{The Web loop is not the native loop}

Fixing ownership prevents use-after-destruction, but it does not make the two schedulers
semantically identical.  At the audited revision, the Web path has five observable differences:

\begin{enumerate}
  \item \texttt{SuppressDraw} is ignored;
  \item variable time-step mode is ignored and the loop always advances as fixed-step;
  \item \texttt{IsRunningSlowly} remains false;
  \item \texttt{EndRun} and \texttt{AfterLoop} are never called;
  \item elapsed time is clamped to a hard-coded 250\,ms rather than the configured
        \texttt{MaxElapsedTime}.
\end{enumerate}

The clock source differs too: the Web route uses \texttt{SDL\_GetTicks}, whereas the native
route uses a performance counter and delay.  None of these differences necessarily prevents
a demo from drawing.  Each can change game behavior, tests, shutdown, or timing-sensitive
network code, so they belong in the compatibility contract rather than in an implementation
footnote.

\section{How to test a host-owned loop}

A useful regression test cannot stop at ``the callback was registered.''  It must keep the
page alive until at least one frame, assert that the game object is still alive inside that
frame, and communicate the result through a browser-visible channel.  It should separately
exercise cancellation and shutdown if those semantics are claimed.

The minimum diagnostic trace is equally small: construction, registration, apparent unwind,
first callback, and destruction.  On a conventional process these events seem obvious; under
a foreign scheduler, spelling them out is often the fastest route to the bug.  The transferable
rule is that a function which syntactically does not return may still run cleanup machinery.
Ownership must follow the host's actual control flow, not the C\texttt{++} surface illusion.
